\chapter{Bessel関数}
    \section{定義}
        $x\in\realNumbers,\;z\in\complexNumbers$に対して$\NapierE^{iz\sin x}$のFourier 級数展開により「$n$次の第一種 Bessel 関数$J_n(z)$」を次のように定義する。
        \begin{equation}
            \label{eqn:implicit definition of 1st kind n-th Bessel function}
            \NapierE^{iz\sin x} = \sum_{n = -\infty}^\infty J_n(z)\NapierE^{inx}
        \end{equation}
        すなわち
        \begin{equation}
            \label{eqn:explicit definition of 1st kind n-th Bessel function}
            J_n(z) = \frac{1}{2\pi}\integrate{-\pi}{\pi}{\NapierE^{iz\sin x}\NapierE^{-inx}}{}{x}
        \end{equation}
        \subsection{簡単化}
            式\eqref{eqn:explicit definition of 1st kind n-th Bessel function}は次のようにして簡単化できる。
            \[ J_n(z) = \frac{1}{2\pi}\integrate{-\pi}{\pi}{\exp i\left(z\sin x -nx\right)}{}{x} = \frac{1}{2\pi}\integrate{-\pi}{\pi}{\cos\left(z\sin x -nx\right)}{}{x} + \frac{1}{2\pi}\integrate{-\pi}{\pi}{i\sin\left(z\sin x -nx\right)}{}{x} \]
            第1項は偶関数の対称積分であり，第 2 項は奇関数の対称積分なので0だから結局次のようになる。
            \begin{equation}
                \label{eqn:simplified definition of 1st kind n-th Bessel function}
                J_n(z) = \frac{1}{2\pi}\integrate{-\pi}{\pi}{\cos\left(z\sin x -nx\right)}{}{x} = \frac{1}{\pi}\integrate{0}{\pi}{\cos\left(z\sin x -nx\right)}{}{x}
            \end{equation}
    \section{性質}
        \subsection{\texorpdfstring{$J_n(-z) = J_{-n}(z)$}{y 軸に関する鏡像反転}}
            式\eqref{eqn:simplified definition of 1st kind n-th Bessel function}より容易に示せる。
        \subsection{\texorpdfstring{$J_{2m}(-z) = J_{2m}(z),\;J_{2m+1}(-z) = -J_{2m+1}(z)$}{次数と偶関数,奇関数}}
            $m\in\integers$とするとき，次式が成り立つ。
            \[ J_{2m}(-z) = J_{2m}(z),\;J_{2m+1}(-z) = -J_{2m+1}(z) \]
            \begin{proof}
                \quad\par
                $n=2m$とする。
                \begin{align*}
                    J_{2m}(-z) &= \frac{1}{2\pi}\integrate{-\pi}{\pi}{\cos(-z\sin x - 2mx)}{}{x} \\
                    &= \frac{1}{2\pi}\integrate{-2\pi}{0}{\cos(z\sin y - 2m(y+\pi))}{}{y} \quad (\text{変数変換 }x=y+\pi) \\
                    &= \frac{1}{2\pi}\integrate{-\pi}{\pi}{\cos(z\sin y - 2m(y+\pi))}{}{y} = \frac{1}{2\pi}\integrate{-\pi}{\pi}{\cos(z\sin y - 2my)}{}{y} = J_{2m}(z)
                \end{align*}
                同様にして$J_{2m+1}(-z) = -J_{2m+1}(z)$も示せる。
            \end{proof}
    \section{Maclaurin 展開}
        式\eqref{eqn:implicit definition of 1st kind n-th Bessel function}の左辺を指数関数に関して Maclaurin 展開すると$x$を変数とした Fourier 級数展開が得られることを利用して$J_n(z)$の Maclaurin 展開を求める。
        \[ \NapierE^{iz\sin x} = \exp\left( z\frac{\NapierE^{ix} - \NapierE^{-ix}}{2} \right) = \sum_{k=0}^\infty \frac{1}{k!} \left(\frac{z}{2}\right)^k (\NapierE^{ix}-\NapierE^{-ix})^k \] %= \sum_{k=0}^\infty \frac{1}{k!} \left(\frac{z}{2}\right)^k \sum_{l=0}^k \binom{k}{l} (-1)^{k-l} \NapierE^{i(2l-k)x}
        $n\geq 0$のとき，$(\NapierE^{ix}-\NapierE^{-ix})^k$の展開で$\NapierE^{inx}$が生じるのは$k=n+2l\;(l=0,1,2,\dotsc)$のときで，その係数は$(-1)^l\binom{k}{l} = (-1)^l\binom{n+2l}{l}$である。
        $n<0$のとき，$(\NapierE^{ix}-\NapierE^{-ix})^k$の展開で$\NapierE^{inx}$が生じるのは$k=-n+2l\;(l=0,1,2,\dotsc)$のときで，その係数は$(-1)^{k-l}\binom{k}{l} = (-1)^{-n+l}\binom{-n+2l}{l}$である。
        式\eqref{eqn:implicit definition of 1st kind n-th Bessel function}の両辺の Fourier 展開基底の係数同士が一致せねばならぬから
        \[
            J_n(z) =
            \begin{dcases}
                \sum_{l=0}^\infty \frac{(-1)^l}{(n+2l)!}\left(\frac{z}{2}\right)^{n+2l}\binom{n+2l}{l} = \sum_{l=0}^\infty \frac{(-1)^l}{l!(n+l)!}\left(\frac{z}{2}\right)^{n+2l} & (n\geq 0) \\
                \sum_{l=0}^\infty \frac{(-1)^{-n+l}}{(-n+2l)!}\left(\frac{z}{2}\right)^{-n+2l}\binom{-n+2l}{l} = \sum_{l=0}^\infty \frac{(-1)^{-n+l}}{l!(-n+l)!}\left(\frac{z}{2}\right)^{-n+2l} & (n < 0)
            \end{dcases}
        \]